\documentclass{article}
\usepackage[utf8]{inputenc}
\usepackage[russian]{babel}

\title{Пример научной статьи}
\author{Автор}
\date{\today}

\begin{document}

\maketitle

\section{Введение}
В этой статье мы рассмотрим основные принципы написания научных статей в формате \LaTeX. \LaTeX{} является мощным инструментом для создания профессионально выглядящих документов, особенно при работе с математическими формулами и библиографией. Введение обычно содержит обзор темы, цели исследования и краткое изложение основных результатов.

\section{Основная часть}
Основная часть статьи содержит детальное описание методов, результатов и анализа. В этой части обычно приводятся теоретические выкладки, экспериментальные данные и их интерпретация. Для математических выражений используются различные окружения, такие как equation, align и другие.

\section{Заключение}
В заключении подводятся итоги исследования, формулируются основные выводы и указываются направления для дальнейшей работы. Заключение играет важную роль в восприятии статьи, так как именно здесь читатель получает краткое резюме проделанной работы и понимает значимость полученных результатов.

\bibliographystyle{plain}
\bibliography{references}

\end{document}